
Die inverse Kinematik berechnet einen vorgebenen Pfad des Endeffektors die benötigten Winkelstellungen der drei Motoren des Deltaroboters. 
Für jeden Motorarm wird dabei separat bestimmt, unter welchem Winkel sich der obere Arm befinden muss, damit der zugehörige Punkt am Endeffektor erreicht werden kann.

Der Zielpunkt des Endeffektors wird im Programm mit dem Punkt \(A\) beschrieben. 
Da jeder Motor an einem anderen Verbindungspunkt des Endeffektors angreift, wird für jeden Motor zunächst ein Offset berücksichtigt:

\[
A_i = A + \Delta A_i
\]

Dabei ist \(A_i\) der Zielpunkt für den jeweiligen Motor \(i\), \(A\) der Mittelpunkt des Endeffektors und \(\Delta A_i\) der motorabhängige Versatz.

Die Position des Motors wird mit \(B\) bezeichnet. Zusätzlich besitzt jeder Motor eine Drehachse \(n_B\). 
Aus der Motorgeometrie ergeben sich damit:

\[
B =
\begin{bmatrix}
x_{B} \\ y_{B} \\ z_{B}
\end{bmatrix},
\qquad
n_{B} =
\frac{1}{\lVert n_{B} \rVert}
\begin{bmatrix}
n_x \\ n_y \\ n_z
\end{bmatrix}
\]

Der obere Arm besitzt die Länge \(l_o\), der untere Arm die Länge \(l_u\). 
Der gesuchte Ellbogenpunkt \(E_i\) muss daher zwei Bedingungen erfüllen:

\[
\lVert E - B \rVert = l_o
\]

\[
\lVert E - A_i \rVert = l_u
\]

Geometrisch entspricht dies dem Verschnitt eines Kreises mit einer Kugel. 
Der Kreis liegt hierbei um den Motorpunkt \(B\) mit dem Radius \(l_o\) und dem Normalenvektor \(n_{B}\), die zweite Kugel liegt um den Zielpunkt \(A_i\) mit dem Radius \(l_u\). 
Es ergeben sich zwei Schnittpunkte wo der Ellbogenpunkt liegen kann.

\[
E_1, \quad E_2
\]

Da der Deltaroboter nach unten arbeitet, wird der Punkt mit der kleineren \(z\)-Koordinate ausgewählt:

\[
E =
\begin{cases}
E_1, & \text{falls } z_{E_1} < z_{E_2} \\
E_2, & \text{sonst}
\end{cases}
\]

Damit ist die Position des Ellbogens vollständig bestimmt. 
Aus dem Ellbogenpunkt kann nun der Richtungsvektor des oberen Arms berechnet werden:

\[
\vec{r}_o = E - B
\]

\[
\hat{r}_o = \frac{\vec{r}_o}{\lVert \vec{r}_o \rVert}
\]

Ebenso ergibt sich der Richtungsvektor des unteren Arms aus:

\[
\vec{r}_u = E - A_i
\]

\[
\hat{r}_u = \frac{\vec{r}_u}{\lVert \vec{r}_u \rVert}
\]

Zur Berechnung des Motorwinkels wird eine lokale Bezugsebene am Motor definiert. 
Die Nullrichtung zeigt vom Roboterzentrum \(C\) zum Motorpunkt \(B\):

\[
\vec{r}_0 = B - C
\]

Da diese Richtung senkrecht zur Motorachse liegen muss, wird sie zunächst auf die Motorebene projiziert:

\[
\vec{r}_{0,\perp}
=
\vec{r}_0
-
(\vec{r}_0 \cdot n_B)n_B
\]

Anschließend wird die Richtung normiert:

\[
\hat{r}_{0,\perp}
=
\frac{\vec{r}_{0,\perp}}{\lVert \vec{r}_{0,\perp} \rVert}
\]

Eine zweite lokale Achse wird über das Kreuzprodukt gebildet:

\[
\hat{r}_{90}
=
\frac{n_B \times \hat{r}_{0,\perp}}
{\lVert n_B \times \hat{r}_{0,\perp} \rVert}
\]

Der Vektor vom Motorzentrum zum Ellbogenpunkt lautet:

\[
\vec{r}_E = E - B
\]

Dieser Vektor wird auf die beiden lokalen Richtungen projiziert:

\[
x_\text{lokal} = \vec{r}_E \cdot \hat{r}_{0,\perp}
\]

\[
y_\text{lokal} = \vec{r}_E \cdot \hat{r}_{90}
\]

Der Motorwinkel ergibt sich anschließend mit der Arkustangens-Funktion:

\[
\vartheta = \operatorname{atan2}(y_\text{lokal}, x_\text{lokal})
\]

Dieser Ablauf wird für alle drei Motoren durchgeführt. 
Das Ergebnis der inversen Kinematik ist somit eine Liste der drei Motorwinkel.
Falls für einen Punkt kein gültiger Schnitt der geometrischen Bedingungen gefunden wird, wird dieser Punkt als nicht erreichbar markiert. 
Damit kann geprüft werden, ob eine vorgegebene Bahn vollständig innerhalb des Arbeitsraums des Roboters liegt.


